\chapter[Approximating LLMs by WFAs]{Approximating Language Models by Weighted Finite Automata}
\label{ch12:top}

Chapter~11 did \emph{not} prove that transformers are weighted finite automata. It suggested something more modest: in the single-pass regime, finite-state approximations are plausible enough to study seriously. This chapter turns that suggestion into a program.

The most important question is not ``Which extraction algorithm is best?'' It is:
\begin{quote}
What exactly do we want the WFA to approximate?
\end{quote}
Without a precise answer, one can slide too quickly between four different things:
\begin{itemize}
  \item exact structure of the original LLM,
  \item approximation on finite data,
  \item exact rationality of the surrogate,
  \item asymptotic conclusions about the original model.
\end{itemize}
This chapter tries to keep those levels separate.

\section{Approximation Targets}

A WFA surrogate might be asked to preserve any of the following.
\begin{enumerate}
  \item \textbf{String probabilities on a bounded set.} Approximate $\mu(w)$ well for strings $w$ up to some maximum length.
  \item \textbf{Next-token conditionals on bounded contexts.} Approximate $p(\sigma \mid h)$ well for histories $h$ in some finite family.
  \item \textbf{Length statistics.} Approximate the length PGF
  \[
    P(z)=\sum_{n \ge 0} p_n z^n.
  \]
  \item \textbf{Support structure.} Approximate which strings are effectively possible or impossible.
  \item \textbf{Low-rank Hankel structure.} Approximate the string function by one whose Hankel matrix has small rank.
\end{enumerate}

These targets are related, but not equivalent. In particular, good approximation on finite lengths or on a validation corpus does \emph{not} automatically imply that the surrogate and the teacher share the same dominant singularities or the same asymptotic regime.

\section{Line 1: Black-Box Distillation}

The simplest idea is to treat the language model as a teacher and fit a small WFA student.

There are two distinct versions of this idea.

\paragraph{Score fitting.}
One may fit a WFA-style string scorer to match values such as
\[
  s(w)=\log \mu(w)
\]
on a finite corpus. This is useful if the downstream task only needs relative ranking or rescoring. But a fitted score is not automatically a normalized probability distribution.

\paragraph{Distribution fitting.}
If the goal is to minimize KL divergence or total variation against a genuine probability model, then the surrogate must itself be a proper probabilistic object. In the WFA setting that means using a nonnegative, properly normalized automaton rather than an arbitrary real-valued string function.

The advantage of black-box distillation is practical flexibility: it can be applied even when the internal architecture of the teacher is unavailable. The limitation is equally clear: success on a dataset is only evidence about the chosen target and the sampled region of string space.

\section{Line 2: Structural Equivalences}

Some neural architectures are exactly finite-state or WFA-like in a much stronger sense.

\begin{theorem}[Structural equivalence, black-box family of results]
Certain recurrent architectures can be written in a form where the update on input symbol $x_t$ is literally of the form
\[
  h_t = A_{x_t} \odot h_{t-1}
\]
for an appropriate semiring operation $\odot$, making the model a rational recurrence / weighted-automaton-style computation.
\end{theorem}

This is the cleanest regime analytically because the equivalence is exact, not approximate. In that case Chapter~7 applies directly.

The reason this line matters for the present book is not that transformers reduce to these models. It is that exact finite-state structure really does exist inside the broader neural-model landscape, so the rational viewpoint is not merely decorative.

\section{Line 3: Query Learning}

A different strategy is to ask the model targeted questions rather than fit it passively on a fixed corpus.

Angluin's classical $L^*$ algorithm learns a minimal DFA by using membership and counterexample queries. For language models, that picture is only partially transferable:
\begin{itemize}
  \item it is natural for accept/reject or thresholded support questions,
  \item it is much less straightforward for real-valued probabilities,
  \item and exact equivalence queries are essentially unavailable in practice.
\end{itemize}

Still, the idea of \emph{adaptive querying} is valuable. It forces the learner to seek out structurally informative strings instead of hoping that a passive corpus already contains them.

\section{Line 4: Spectral and Hankel Methods}

This is the mathematically richest line.

\begin{definition}[Hankel matrix]
Given a string function $f : \Sigma^* \to \mathbb{R}$, its Hankel matrix is the bi-infinite matrix
\[
  H_f(u,v) = f(uv),
\]
indexed by prefixes $u$ and suffixes $v$.
\end{definition}

The importance of the Hankel matrix is captured by a fundamental theorem.

\begin{theorem}[Finite Hankel rank, black box]
A string function is realizable by a finite-state weighted automaton if and only if its Hankel matrix has finite rank. The minimal number of WFA states equals that rank.
\end{theorem}

This is the weighted analogue of the Myhill--Nerode theorem. In the Boolean case, finite-state structure means finitely many future behaviors. In the weighted case, it means the future-behavior rows span a finite-dimensional linear space.

Spectral learning takes a large finite block of the Hankel matrix, computes a low-rank factorization, and reads off a small WFA from that factorization. The exact formulas depend on the chosen finite prefix/suffix bases and on pseudoinverses, so we do not reproduce them in detail here. The core idea is enough: small Hankel rank means small WFA.

In one-letter settings, operator-theoretic results such as the Adamyan--Arov--Krein theorem say that the best low-rank approximation can be taken \emph{within} the Hankel class, not just among arbitrary low-rank matrices. This is mathematically beautiful, but also far beyond the level of a first pass through the subject. For our purposes, it is enough to know that one-letter Hankel approximation is the cleanest controlled setting.

\section{One-Letter Projections: Useful but Limited}

Collapsing a model to a one-letter alphabet means forgetting token identities and keeping only length. Concretely, if
\[
  p_n = \Pr(|W|=n),
\]
then the one-letter projection is exactly the length PGF
\[
  P(z) = \sum_{n \ge 0} p_n z^n.
\]
This is extremely useful for studying length tails, stopping behavior, and the singularity structure of the length distribution.

But it is \emph{not} the same as the counting-type typical-set generating function of Chapter~9. In particular, the one-letter projection is the right object for length statistics, not the object whose radius directly encodes Shannon entropy via
\[
  h = \log_2(1/R).
\]
So the one-letter projection is a feature, not a bug --- but it answers a specific question, not every question.

\section{Scale Realities}

All four approximation lines become much harder at modern LLM scale.

\begin{itemize}
  \item Distillation is limited by how much of string space the training sample reveals.
  \item Query learning is limited by the impossibility of exact equivalence checks.
  \item Spectral methods are limited by very large effective Hankel rank and by sampling noise.
  \item Exact structural equivalences usually apply only to restricted model classes.
\end{itemize}

This is why small surrogates and controlled synthetic settings matter. For example, one can begin with models trained on data generated from a known PCFG or from a known finite-state source. In those cases the underlying formal structure is known in advance, which makes approximation error interpretable.

\section{What a WFA Surrogate Actually Buys}

Suppose we construct a WFA surrogate
\[
  \widehat M = (\hat\alpha, \{\hat A_\sigma\}, \hat\beta)
\]
with scalar length generating function
\[
  \widehat F(z) = \hat\alpha (I-z\hat A)^{-1} \hat\beta,
  \qquad
  \hat A = \sum_\sigma \hat A_\sigma.
\]
Then two things are true.

\begin{enumerate}
  \item The surrogate has an exact rational length generating function, so Chapters~4 and~7 determine its pole-based asymptotics.
  \item Those conclusions apply \emph{exactly} to the surrogate and only \emph{indirectly} to the original LLM.
\end{enumerate}

So the rational toolkit becomes available, but only for the object we have actually extracted. Whether the surrogate's poles are informative about the teacher depends entirely on the chosen approximation target and on how well that target was matched.

A safe formulation is this: WFA approximation is a way of \emph{proposing} a rational model for some aspect of an LLM. It is not, by itself, a proof that the original LLM has rational asymptotics.

\section*{Preview}

The next chapters use this chapter's distinction aggressively.
\begin{itemize}
  \item Chapter~13 and beyond will study exact analytic objects such as partition functions and length PGFs.
  \item Later empirical chapters will ask whether simple rational or algebraic surrogates actually fit observed data.
\end{itemize}
The approximation target chosen here will determine what those later claims can legitimately mean.

\section*{Further Reading}

\begin{itemize}
  \item \textbf{G. Lecorvé and P. Motlicek}, ``Conversion of recurrent neural network language models to weighted finite-state transducers'' (2012) \cite{LecorveMotlicek2012} --- an early example of black-box distillation.
  \item \textbf{A. T. Suresh et al.}, ``Approximating probabilistic models as weighted finite automata'' (2021) \cite{SureshEtAl2021} --- a modern distillation reference.
  \item \textbf{G. Weiss, Y. Goldberg, and E. Yahav}, ``Extracting automata from recurrent neural networks using queries and counterexamples'' (2018) \cite{WeissGoldbergYahav2018} --- for the query-learning perspective.
  \item \textbf{G. Rabusseau, T. Li, and D. Precup}, ``Connecting weighted automata and recurrent neural networks through spectral learning'' (2019) \cite{RabusseauLiPrecup2019} --- for spectral reconstruction ideas.
  \item \textbf{C. Lacroce, P. Panangaden, and G. Rabusseau}, ``Extracting weighted finite automata from recurrent neural networks using constrained AAK theory'' (2021) \cite{LacrocePanangadenRabusseau2021} --- for the one-letter AAK perspective.
  \item \textbf{J. Berstel and C. Reutenauer}, \emph{Noncommutative Rational Series with Applications} (2011) \cite{BerstelReutenauer2011} --- the standard algebraic reference on rational series and Hankel rank.
\end{itemize}

\begin{exercise}
For the Fibonacci WFA from Chapter~7, compute
\[
  F(z)=\alpha(I-zA_a)^{-1}\beta
\]
explicitly and identify its poles.
\end{exercise}

\begin{exercise}
Consider the scalar string function over the one-letter alphabet defined by
\[
  f(1^n)=n+1.
\]
Show that its Hankel matrix has rank $2$ and exhibit a 2-state WFA realizing it.
\end{exercise}
